\subsection{Phase Tower Geometry Interface: Principal \texorpdfstring{$U(1)$}{U(1)} Bundle, Area Weight, and Precision Amplification Law}\label{subsec:unit-disk-phase-tower}
The ``phase compactification gate'' of this appendix can be placed within a more general geometric interface: each ``lift'' structurally manifests as a principal \(U(1)\) bundle (phase fiber), while ``near-zero weighting'' is automatically produced by the Jacobian pullback of the uniform measure.
This subsection provides only the minimal formalism that can be directly interfaced with the gate language of this paper.

\begin{remark}[Standard chain: projectivization \(\to\) Hopf principal \texorpdfstring{$U(1)$}{U(1)} bundle \(\to\) stereographic volume-form pullback]\label{rem:phase-tower-standard-chain}
Projectivization and the Hopf fibration strictly realize ``one lift adds one global phase fiber'' as a principal \(U(1)\) bundle; pullback of the volume form under stereographic projection yields the \((1+|x|^{2})^{-n}\)-type weight decay. These are all standard geometric facts; see \cite{Husemoller1994FiberBundles}.
\end{remark}

\paragraph{Phase tower: principal \texorpdfstring{$U(1)$}{U(1)} bundle and Hopf fibration}
The Hopf fibration provides the standard example
\[
S^{1}\hookrightarrow S^{3}\twoheadrightarrow S^{2},
\]
and for any \(n\ge 1\) gives
\[
S^{1}\hookrightarrow S^{2n+1}\twoheadrightarrow \mathbb{CP}^{n}.
\]
This formally expresses the structural primitive ``one lift adds one global phase fiber'' (see \cite{Husemoller1994FiberBundles}).
In the context of this paper, the phase gates provided by \(\upsilon\) and \(\iota\) can be viewed as the minimal instance of compressing an unbounded coordinate onto a compact phase fiber.

\begin{proposition}[Projectivization and principal \texorpdfstring{$U(1)$}{U(1)} bundle (phase tower primitive)]\label{prop:phase-tower-principal-u1}
For any \(n\ge 1\), let
\(
S^{2n+1}=\{z\in\CC^{n+1}:\|z\|=1\}
\), and define the projection
\[
p:S^{2n+1}\to \mathbb{CP}^{n},\qquad p(z):=[z]
\]
(\([z]\) denotes the complex line spanned by \(z\)).
Then \(U(1)\) acts freely on \(S^{2n+1}\) via \(e^{\mathrm{i}\phi}\cdot z:=e^{\mathrm{i}\phi}z\), and \(p\) is the quotient map; therefore \(p\) constitutes a principal \(U(1)\) bundle.
When \(n=1\), \(p:S^{3}\twoheadrightarrow\mathbb{CP}^{1}\cong S^{2}\) is the Hopf fibration.
\end{proposition}
\begin{proof}
The action of the phase group \(U(1)\) preserves the norm, and \(e^{\mathrm{i}\phi}z=z\Rightarrow e^{\mathrm{i}\phi}=1\) (freeness).
Two points \(z,z'\in S^{2n+1}\) have the same projective image if and only if \(z'=\lambda z\) (\(\lambda\in\CC^\times\)); from \(\|z\|=\|z'\|=1\) one obtains \(|\lambda|=1\), so \(\lambda\in U(1)\).
Therefore the fibers of \(p\) are exactly the \(U(1)\)-orbits, and \(p\) is the standard model of a principal \(U(1)\) bundle (see \cite{Husemoller1994FiberBundles}). \qed
\end{proof}

\paragraph{Local phase and global phase: projectivization and \texorpdfstring{$\det:U(n)\to U(1)$}{det}}
When encoding multi-axis information as a complex vector \(z\in\CC^{n+1}\), one may write the polar decomposition \(z_k=r_k e^{\mathrm{i}\theta_k}\) for each coordinate, so each axis carries a \textbf{local phase} \(\theta_k\).
At the same time there exists a common \textbf{global phase} action \(z\mapsto e^{\mathrm{i}\phi}z\) (\(\phi\in\RR\)), whose quotient space is the complex projective space:
\[
\mathbb{CP}^{n}\ \cong\ (\CC^{n+1}\setminus\{0\})/U(1).
\]
Therefore, the rigidified version of ``after merging there is only one angle'' can be stated as: the global phase is the fiber coordinate of the principal \(U(1)\) bundle, while the relative structure resides on the base space \(\mathbb{CP}^{n}\).
At the matrix level, the determinant \(\det(U)\in U(1)\) of any \(U\in U(n)\) provides a unique ``total phase readout,'' with kernel \(SU(n)\); therefore \(\det:U(n)\to U(1)\) provides an algebraic global phase extraction caliber.

\paragraph{Area weight and Jacobian: structural origin of near-zero weighting}
In the one-dimensional case, the derivative of \(\upsilon(x)=\frac{1}{\pi}\arctan(x)\) yields the natural weight \(\dd\mu(x)=\frac{1}{\pi(1+x^2)}\dd x\) (see \eqref{eq:unit-circle-phase-derivative}--\eqref{eq:unit-circle-phase-cauchy-measure}).
This weight is not an externally imposed rule but rather the Jacobian of the compactification change of variables.
In two or higher dimensions, the analogous phenomenon is given by ``the density factor of the uniform area element/volume element in projective coordinates.''
For example, in the stereographic projection coordinates \(z=X+\mathrm{i}Y\in\CC\) on the Riemann sphere \(S^{2}\), the round metric satisfies
\[
ds^{2}=\frac{4}{(1+|z|^{2})^{2}}\,(\dd X^{2}+\dd Y^{2})
=\frac{4}{(1+|z|^{2})^{2}}\,|\,\dd z\,|^{2},
\]
so the pullback density of the uniform spherical area element is
\[
\boxed{\ \dd A_{S^{2}}=\frac{4}{(1+|z|^{2})^{2}}\,\dd A_{\RR^{2}}\ }.
\]
Therefore, in planar coordinates, larger \(|z|\) corresponds to smaller spherical area; ``near-zero weighting'' is automatically produced by the Jacobian rather than imposed by an additional rule.

\begin{proposition}[Volume-form pullback under stereographic projection (general dimension)]\label{prop:stereographic-volume-form}
Let \(\sigma:S^{n}\setminus\{N\}\to\RR^{n}\) be the stereographic projection from the north pole \(N\), with \(x\in\RR^{n}\) as image-space coordinates.
Let \(\sigma^{-1}:\RR^{n}\to S^{n}\setminus\{N\}\) be its inverse map.
Then the pullback of the spherical metric under \(\sigma^{-1}\) satisfies the conformal scaling
\[
(\sigma^{-1})^{\ast}g_{S^{n}}=\left(\frac{2}{1+\|x\|^{2}}\right)^{2}g_{\mathrm{Euc}},
\]
so the pullback of the spherical volume form satisfies
\begin{equation}\label{eq:stereographic-volume-general}
\boxed{\;
(\sigma^{-1})^{\ast}\dd V_{S^{n}}
 =\left(\frac{2}{1+\|x\|^{2}}\right)^{n}\dd V_{\RR^{n}}
=\frac{2^{n}}{(1+\|x\|^{2})^{n}}\,\dd x_{1}\cdots \dd x_{n}\; }.
\end{equation}
In particular, when \(n=1\) this density differs from \(\dd\mu(x)=\frac{1}{\pi(1+x^{2})}\dd x\) only by a normalization constant, while for \(n=2\) it degenerates to the expression \(\frac{4}{(1+\|x\|^{2})^{2}}\) above.
\end{proposition}
\begin{proof}
Stereographic projection provides the standard conformal coordinates of the sphere on \(\RR^{n}\); the conformal factor is \(\Omega(x)=\frac{2}{1+\|x\|^{2}}\), i.e., \((\sigma^{-1})^{\ast}g_{S^{n}}=\Omega(x)^{2}g_{\mathrm{Euc}}\).
Under conformal scaling, the volume form scales by \(\Omega^{n}\), and therefore
\(
(\sigma^{-1})^{\ast}\dd V_{S^{n}}=\Omega(x)^{n}\dd V_{\RR^{n}}
\), yielding \eqref{eq:stereographic-volume-general}. \qed
\end{proof}

\paragraph{Precision amplification law: dimensional reduction by compactification pushes dynamic range into least-significant-bit precision}
From \(\upsilon^{-1}(u)=\tan(\pi u)\) (equation~\eqref{eq:unit-circle-phase-inverse}) one obtains the error amplification
\begin{equation}\label{eq:app-precision-amplification}
\boxed{\ \frac{\dd x}{\dd u}=\pi(1+x^2)\ }.
\end{equation}
Therefore, when \(|x|\) is large, a tiny phase error \(\Delta u\) is amplified to a readout error \(\Delta x\approx \pi x^2\Delta u\) of enormous magnitude; in other words, the dimensional reduction/compactification forcibly transports ``large-number information far from zero'' into the least-significant-bit precision of \(u\).
Roughly, if the quantization precision of \(u\) is \(\Delta u\), then the maximum resolvable magnitude satisfies
\[
|x|_{\max}\approx \cot(\pi\Delta u)\approx \frac{1}{\pi\Delta u}\qquad(\Delta u\to 0).
\]
This provides a verifiable analytic scaling law for the gate-language assertion that ``information is compressed onto precision.''
